\documentclass[11pt]{article}
\usepackage[a4paper,margin=27mm]{geometry}
\usepackage[T1]{fontenc}
\usepackage{lmodern,microtype,amsmath,amssymb,amsthm,mathtools,enumitem}
\usepackage[hidelinks]{hyperref}
\newtheorem{theorem}{Theorem}
\newtheorem{lemma}[theorem]{Lemma}
\newtheorem{corollary}[theorem]{Corollary}
\newtheorem{proposition}[theorem]{Proposition}
\theoremstyle{remark}
\newtheorem{remark}[theorem]{Remark}
\newcommand{\R}{\mathbb R}
\newcommand{\one}{\mathbf 1}
\newcommand{\supp}{\operatorname{supp}}
\newcommand{\sd}{\mathbin{\triangle}}
\newcommand{\dH}{d_{\mathrm H}}
\newcommand{\Dscr}{\mathcal D}
\newcommand{\Escr}{\mathcal E}
\newcommand{\lpr}[1]{\underline{P_{#1}}}
\newcommand{\upr}[1]{\overline{P_{#1}}}
\newcommand{\lowprob}[1]{\underline{p_{#1}}}
\newcommand{\upprob}[1]{\overline{p_{#1}}}
\title{Accuracy Rankings for Every Coherent Set of Desirable Gambles}
\author{}
\date{8 October 2026}
\begin{document}
\maketitle

\noindent
Given any coherent set of desirable gambles, we construct an accuracy ranking
under which it is undominated. The ranking evaluates every measurable
acceptance set and strictly rewards reductions in errors, up to null sets.
The resulting class of rankings satisfies \textbf{(E)}, \textbf{(LA)},
\textbf{(A)}, and \textbf{(S)}, defined below. We also construct bounded
real-valued losses with the same properties.

The proof has three steps. First place the target in a continuous family of
coherent sets. Next use their lower and upper probabilities to construct
rankings that keep every family member undominated. Finally repeat the
construction for each target. The numerical loss uses the same families,
with an additional estimate controlling how much an error reduction lowers loss.

\section{Models, accuracy, and comparisons across worlds}

Fix worlds $\Omega=\{\omega_1,\dots,\omega_n\}$, with $n\ge2$, and the simplex
$\Delta=\{p\in\R^n:p\ge0,\ \sum_i p_i=1\}$.
A gamble $g\in\R^n$ pays $g_i$ at $\omega_i$. A model, or acceptance set,
is a Lebesgue-measurable set $\Dscr\subseteq\R^n$: accept exactly its members. Fix a probability measure
$\mu$ with the same null sets as Lebesgue measure, and write
$\Dscr\approx \Escr$ when $\mu(\Dscr\sd \Escr)=0$, where $\sd$ denotes symmetric difference.
All background comparisons below are understood modulo null sets.

Write $[n]=\{1,\dots,n\}$, $\one=(1,\dots,1)$, and let $e_i$ be the $i$th
unit vector, also written $\one_{\omega_i}$ as an indicator gamble.
Write $p_i=p(\omega_i)$. For $p\in\Delta$, put $\supp p=\{i:p_i>0\}$; for
$S\subseteq[n]$, put $p(S)=\sum_{i\in S}p_i$,
$S^c=[n]\setminus S$, and $\Delta_S=\{p\in\Delta:\supp p\subseteq S\}$.
We use the usual norms $\|x\|_1=\sum_i|x_i|$ and
$\|x\|_\infty=\max_i|x_i|$. The Hausdorff distance $\dH(K,H)$ is measured in
$\|\cdot\|_1$: it is the least distance within which each point of either set
has a point of the other.

A coherent set of desirable gambles excludes $0$, contains every nonzero
$g\ge0$, and is closed under addition and multiplication by positive scalars.
A \emph{credal set} is a nonempty closed convex set $K\subseteq\Delta$ of
probability distributions. Its coherent lower prevision and upper prevision are
\[
 \lpr{K}(g)\coloneqq\min_{m\in K}m\cdot g,\qquad
 \upr{K}(g)\coloneqq\max_{m\in K}m\cdot g=-\lpr{K}(-g).
\]
The lower and upper probabilities of a world are
\[
 \begin{aligned}
 \lowprob{K}(\omega_i)&\coloneqq\lpr{K}(\one_{\omega_i})=\min_{m\in K}m_i,\\
 \upprob{K}(\omega_i)&\coloneqq\upr{K}(\one_{\omega_i})=\max_{m\in K}m_i.
 \end{aligned}
\]
We identify these functions with their vectors of values on the $n$ worlds.
Thus $\lowprob{K}$ and $\upprob{K}$ need not themselves sum to one.
As in the CPM notation, distinguish the strictly positive set from its coherent
completion:
\[
 D_K\coloneqq\{g:\lpr{K}(g)>0\},\qquad
 \Dscr_K\coloneqq D_K\cup\{g:g\ge0,\ g\ne0\}.
\]
For a single distribution $p$, write $D_p=D_{\{p\}}=\{g:p\cdot g>0\}$.
For general acceptance sets we use $\Dscr,\Escr$; they need not be coherent.
The added nonnegative gambles ensure coherence at boundary credal sets and,
as shown below, change $D_K$ only on a null set.

We first record the elementary properties needed to compare these sets and
take limits. They follow directly from minimizing expectations over $K$.
For any credal sets $K,H$ and gambles $g,h$,
\begin{equation}
 \begin{aligned}
 |\lpr{K}(g)-\lpr{K}(h)|&\le\|g-h\|_\infty,\\
 |\lpr{K}(g)-\lpr{H}(g)|&\le\dH(K,H)\|g\|_\infty.
 \end{aligned}\label{eq:phi-estimates}
\end{equation}
The first bound follows by taking minima in
$|m\cdot(g-h)|\le\|g-h\|_\infty$ for $m\in\Delta$.
For the second, choose $m\in K$ minimizing $m\cdot g$ and $m'\in H$ with
$\|m-m'\|_1\le\dH(K,H)$. Then
$\lpr{H}(g)\le m'\cdot g\le\lpr{K}(g)+\dH(K,H)\|g\|_\infty$;
exchange $K,H$ for the reverse bound. Adding a constant shifts every expectation by that constant, since
$m\cdot\one=1$ for every $m\in\Delta$.

The set $\Dscr_{K}$ excludes zero and contains every nonzero nonnegative gamble by
definition. Positive homogeneity of $\lpr{K}$ gives closure under positive
scaling. For addition, use
$\lpr{K}(g+h)\ge\lpr{K}(g)+\lpr{K}(h)$. If one of $g,h\in \Dscr_{K}$ has
strictly positive lower prevision, so does their sum, since the other has
nonnegative lower prevision. Otherwise both are nonzero and nonnegative,
and so is their sum. Thus $\Dscr_{K}$ is coherent.

The lower prevision is therefore continuous, and adding a constant gives
\begin{equation}
 \lpr{K}(g+t\one)=\lpr{K}(g)+t.\label{eq:constant}
\end{equation}
Thus its zero set meets each line parallel to $\one$ exactly once, and is null
by Fubini's theorem. Since
$\{\lpr{K}>0\}\subseteq \Dscr_{K}\subseteq\{\lpr{K}\ge0\}$, we have
$\Dscr_{K}\approx\{\lpr{K}>0\}$.

The ideal acceptance set at $\omega_i$ is $D_i=\{g:g_i>0\}$.
The errors made by $\Dscr$ there are $E_i(\Dscr)=\Dscr\sd D_i$.
The background accuracy order is
\[
 \Dscr\le_i\Escr\quad\Longleftrightarrow\quad E_i(\Dscr)\subseteq E_i(\Escr)\text{ a.e.}
\]
Its strict part $\Dscr<_i\Escr$ means, additionally,
$\mu(E_i(\Escr)\setminus E_i(\Dscr))>0$. Equivalently, $\Dscr\le_i\Escr$ says that, up to
null sets,
\begin{equation}
 \Escr\cap D_i\subseteq \Dscr,\qquad \Dscr\setminus D_i\subseteq \Escr.\label{eq:acceptance}
\end{equation}
Where accepting is correct, $\Dscr$ accepts at least what $\Escr$ accepts; where
accepting is an error, it accepts at most what $\Escr$ accepts. We will also use
$E_i(\Dscr)\sd E_i(\Escr)=\Dscr\sd \Escr$.

\paragraph{Rankings and the requirements.}
A \emph{ranking} assigns each world $\omega_i$ a reflexive, transitive relation
$\preceq_i$ on its evaluation domain, with $\Dscr\preceq_i\Escr$ meaning that $\Dscr$
is at least as accurate as $\Escr$. Its strict part is $\prec_i$; models tie
when both comparisons hold. The relations need not be complete.
A model $\Dscr$ \emph{weakly dominates} $\Escr$ if $\Dscr\preceq_i\Escr$ at every world;
it \emph{dominates} $\Escr$ if, in addition, $\Dscr\prec_i\Escr$ at some world, and
\emph{strictly dominates} $\Escr$ if this holds at every world. A model is
\emph{undominated} if none dominates it. It is \emph{competitive on}
$S\ne\varnothing$ if for every evaluated model $\Escr$ some $i\in S$ has
$\Dscr\preceq_i\Escr$.

Let $\mathcal C$ be the set of all models agreeing a.e. with a coherent set of
desirable gambles. For each ranking $R$ in a class $\mathfrak R$, let
$\mathcal C_R$ be its evaluated, coherent, undominated models. We establish
these requirements:
\begin{enumerate}[leftmargin=*,itemsep=2pt]
\item[\textbf{(E)}] Every model is evaluated by some ranking in $\mathfrak R$.
\item[\textbf{(LA)}] $\mathcal C_R\ne\varnothing$ for every $R\in\mathfrak R$.
\item[\textbf{(A)}] Every member of $\mathcal C$ belongs to some $\mathcal C_R$.
\item[\textbf{(S)}] Each ranking admits a selection $p\mapsto c_p\in\mathcal C_R$
such that $c_p$ is competitive on $\supp p$ and
$\{p:c_p\preceq_i\Dscr\}$ is closed for every world $i$ and evaluated model $\Dscr$.
\item[\textbf{(C)}] If an evaluated model ties at every world with a member of
$\mathcal C_R$, it also belongs to $\mathcal C_R$.
\end{enumerate}
Every ranking constructed below evaluates every measurable model. The symbol
$\approx$ continues to mean agreement a.e.; it does not by definition mean
that two models tie under a ranking.

The next lemma is the link between accuracy and probability. If one acceptance
set makes no more errors than two reference sets at different worlds, their
lower and upper probabilities must be ordered in the same direction.

\begin{lemma}\label{lem:endpoints}
If $\Dscr\le_i\Dscr_{K}$, $\Dscr\le_j\Dscr_{H}$, and $i\ne j$, then
\begin{equation}
 \lowprob{K}(\omega_i)\le\lowprob{H}(\omega_i),\qquad \upprob{K}(\omega_i)\le\upprob{H}(\omega_i).\label{eq:endpoints}
\end{equation}
\end{lemma}
\begin{proof}
On $W_{ij}=\{g:g_i>0\ge g_j\}$, accepting is correct at $i$ and an error at
$j$. Thus \eqref{eq:acceptance} gives
\[
 \Dscr_{K}\cap W_{ij}\subseteq \Dscr\cap W_{ij}\subseteq \Dscr_{H}\cap W_{ij}
 \quad\text{a.e.}
\]
Exchanging $i,j$ gives $\Dscr_{H}\cap W_{ji}\subseteq \Dscr_{K}$ a.e.

If $\lowprob{K}(\omega_i)>\lowprob{H}(\omega_i)$, choose $t$ strictly between them. Then $0<t<1$, and
$g=e_i-t\one$ satisfies
\[
 g_i>0>g_j,\qquad
 \lpr{K}(g)=\lowprob{K}(\omega_i)-t>0>\lowprob{H}(\omega_i)-t=\lpr{H}(g).
\]
By continuity, these strict inequalities persist on an open neighbourhood,
which lies in $(\Dscr_{K}\cap W_{ij})\setminus \Dscr_{H}$. This contradicts its nullity.
Similarly, if $\upprob{K}(\omega_i)>\upprob{H}(\omega_i)$, choose $t$ between them and use
$g=t\one-e_i$: then $g\in W_{ji}$ and
$\lpr{H}(g)=t-\upprob{H}(\omega_i)>0>t-\upprob{K}(\omega_i)=\lpr{K}(g)$, contradicting the
other inclusion.
\end{proof}

We will use a weighted average of lower and upper probabilities to label the
family. The lemma then turns background accuracy comparisons into inequalities
between those labels.

\begin{corollary}\label{cor:cc}
Suppose a family $p\mapsto K_p$ satisfies, for a fixed $0<\lambda<1$,
\begin{equation}
 (1-\lambda)\lowprob{K_p}(\omega_i)+\lambda\upprob{K_p}(\omega_i)=p_i\quad\text{for all }p,i.\label{eq:index}
\end{equation}
Then it satisfies the following condition on comparisons at different worlds:
\begin{equation*}
 \Dscr\le_i\Dscr_{K_p},\quad \Dscr\le_j\Dscr_{K_q},\quad i\ne j
 \quad\Longrightarrow\quad p_i\le q_i.
 \tag{CC}\label{eq:cc}
\end{equation*}
\end{corollary}
\begin{proof}
Take the weighted average in \eqref{eq:endpoints} with $K=K_p$ and $H=K_q$.
\end{proof}

\section{A family containing the target}

Fix the target credal set $M$; our goal is to make $\Dscr_M$ undominated.
We now build a family containing it whose lower and upper probabilities obey
\eqref{eq:index}. This will supply the comparisons needed to define the ranking.

Let $M\subseteq\Delta$ be any credal set, with lower- and upper-probability vectors $\ell=\lowprob{M}$
and $u=\upprob{M}$. Put $L=\sum_i\ell_i$ and $U=\sum_i u_i$. If $M$ is not a
singleton, then $L<1<U$: equality in either bound would force every member
of $M$ to equal that endpoint vector.

Choose the reference point $r$ where the segment from $\ell$ to $u$ crosses the
hyperplane $\sum_i x_i=1$. Explicitly,
\begin{equation}
 \lambda\coloneqq\frac{1-L}{U-L},\qquad
 r\coloneqq(1-\lambda)\ell+\lambda u,\qquad
 d\coloneqq r-\ell=\lambda(u-\ell).\label{eq:anchor}
\end{equation}
Then $0<\lambda<1$, $r\in\Delta$ and $d\ge0$. The reference point need not belong to
$M$. For a singleton $M=\{m\}$, take $\lambda=1/2$, $r=m$ and $d=0$.

We translate and scale $M$ so that the reference point moves from $r$ to $p$:
\begin{equation}
 \theta(p)\coloneqq\min\left\{1,\ \min_{i:d_i>0}\frac{p_i}{d_i}\right\},
 \qquad
 \boxed{\quad K_p\coloneqq p+\theta(p)(M-r).\quad}\label{eq:family}
\end{equation}
The scale is the largest factor at most one for which the translated copy
stays in the simplex: its lowest $i$th coordinate is $p_i-\theta(p)d_i$.
The minimum over an empty set is $+\infty$, so for singleton targets
$\theta\equiv1$ and $K_p=\{p\}$.

\begin{theorem}[Full family coverage]\label{thm:family}
For every target $M$, the family \eqref{eq:family} has the following properties.
\begin{enumerate}[label=(\alph*)]
\item Each $K_p$ is a credal set, and $K_r=M$.
\item The weighted-average identity \eqref{eq:index}, and hence (CC), holds.
\item $\upprob{K_p}(\omega_i)\le p_i/\lambda$. In particular,
$K_p\subseteq\Delta_{\supp p}$: if the index assigns zero probability
to a world, every distribution in $K_p$ does too.
\item There is a finite constant $A$ such that
$\dH(K_p,K_q)\le A\|p-q\|_1$. In particular,
$\mu(\Dscr_{K_p}\sd \Dscr_{K_q})\to0$ as $q\to p$.
\end{enumerate}
The target may meet any boundary of the simplex, including its vertices.
\end{theorem}
\begin{proof}
Each $K_p$ is a nonempty compact convex set whose points have coordinate sum
one. Since $\theta(p)d_i\le p_i$, its lower-probability vector is
\[
 \lowprob{K_p}=p+\theta(p)(\ell-r)=p-\theta(p)d\ge0.
\]
Thus $K_p\subseteq\Delta$. Also $r_i\ge d_i$, so $\theta(r)=1$ and $K_r=M$.

The upper-probability vector is $\upprob{K_p}=p+\theta(p)(u-r)$. Their weighted
average is therefore
\[
 \begin{aligned}
 (1-\lambda)\lowprob{K_p}+\lambda\upprob{K_p}
 &=p+\theta(p)\big((1-\lambda)\ell+\lambda u-r\big)\\
 &=p.
 \end{aligned}
\]
This proves (b). Since the lower probabilities are nonnegative, the same identity
gives $p_i\ge\lambda\upprob{K_p}(\omega_i)$. If $p_i=0$, every member of $K_p$ must
therefore have zero $i$th coordinate, proving (c).

For (d), $\theta$ is a minimum of finitely many affine functions, hence
Lipschitz. Pairing the images of the same $m\in M$ gives
\begin{equation}
 \dH(K_p,K_q)\le\|p-q\|_1+2|\theta(p)-\theta(q)|,\label{eq:hausdorff}
\end{equation}
since $\|m-r\|_1\le2$. Hausdorff convergence implies pointwise convergence
of the lower previsions by \eqref{eq:phi-estimates}. Their signs therefore
converge outside the null zero set of the limit. Since $\mu$ is finite,
dominated convergence gives the asserted continuity in measure.
\end{proof}

\paragraph{The vacuous target.}
For $M=\Delta$, the reference point is $r=\one/n$ and $\lambda=1/n$. Writing
$m(p)=\min_i p_i$, the family becomes
\begin{equation}
 \boxed{\quad K_p=p-m(p)\one+n\,m(p)\Delta.\quad}\label{eq:vacuous}
\end{equation}
It equals $\Delta$ at the uniform index and $\{p\}$ on the boundary. Its
endpoints are $\lowprob{K_p}(\omega_i)=p_i-m(p)$ and
$\upprob{K_p}(\omega_i)=p_i+(n-1)m(p)$, so
\[
 p_i=\frac{n-1}{n}\lowprob{K_p}(\omega_i)+\frac1n\upprob{K_p}(\omega_i).
\]
This family need not satisfy the earlier upper-probability condition (U$^*$).
Appendix~\ref{app:ustar} explains why that stronger condition excludes some
targets even though (CC) permits full coverage.

\section{Constructing the accuracy ranking}

The family is now fixed and contains the target. We next assign comparison
levels to every acceptance set, then adjust them just enough to make strict
background improvements strict without allowing any family member to be dominated.

Write $C_p=\Dscr_{K_p}$ for a selected family member; in particular,
$C_r=\Dscr_M$. This notation keeps it distinct from the precise-probability
set $D_p$. Call $p$ \emph{available for $\Dscr$ at world $i$} when
$\Dscr\le_iC_p$: $\Dscr$ makes no error there that $C_p$ avoids. Let $\mathcal P_i(\Dscr)$ collect those indices, and $T(\Dscr)$ the worlds
where at least one exists. Define comparison levels $a_i$ and the remaining
allowance $s$ by
\[
 \mathcal P_i(\Dscr)\coloneqq\{p:\Dscr\le_iC_{p}\},\qquad
 T(\Dscr)\coloneqq\{i:\mathcal P_i(\Dscr)\ne\varnothing\},
\]
\[
 a_i(\Dscr)\coloneqq\max_{p\in\mathcal P_i(\Dscr)}p_i\quad(i\in T(\Dscr)),\qquad
 s(\Dscr)\coloneqq1-\sum_{i\in T(\Dscr)}a_i(\Dscr).
\]
Here $a_i$ takes the largest available coordinate, and $s$ is what remains
after adding these values and subtracting from one.
These sets of indices are compact: the function
$p\mapsto\mu(E_i(\Dscr)\setminus E_i(C_{p}))$ is continuous and its zero set is
$\mathcal P_i(\Dscr)$. Thus the maxima exist. Set $\bar a_i(\Dscr)=a_i(\Dscr)$ at worlds with an available index and $\bar a_i(\Dscr)=0$ otherwise, so $s(\Dscr)=1-\sum_i\bar a_i(\Dscr)$.

The next lemma shows that the comparison levels sum to at most one.
Equality identifies exactly the selected family; the amount left over will
allow us to make strict accuracy improvements strict in the final ranking.

\begin{lemma}[Comparison levels and the remaining allowance]\label{lem:levels}
The following hold.
\begin{enumerate}[label=(\alph*)]
\item If $\Dscr\le_i\Escr$ and $\Escr$ has an available index at $i$, then $\Dscr$ has an available index there and
$a_i(\Dscr)\ge a_i(\Escr)$.
\item $a_i(C_{p})=p_i$. If $C_{p}\le_iC_{q}$ and $p_i=q_i$, then $p=q$.
\item $s(\Dscr)\ge0$, with equality exactly when $\Dscr\approx C_{p}$ for some $p$.
\end{enumerate}
\end{lemma}
\begin{proof}
(a) follows from inclusion of the sets of available indices. For (b), $p$ is available for $C_p$ at every world. Any other available index $q$ at $i$, paired with the index $p$ at a different
world in (CC), gives $q_i\le p_i$. If $C_{p}\le_iC_{q}$, pairing $q$ at $i$
with $p$ at each $j\ne i$ instead gives $p_j\le q_j$; equality at $i$ and
equal coordinate sums then force $p=q$.

For (c), let $T=T(\Dscr)\ne\varnothing$, and choose $p^j$ attaining the maximum in $a_j(\Dscr)$
for each $j\in T$. By (CC), $p^j_k\ge p^k_k=\bar a_k(\Dscr)$ for $k\in T$,
with equality when $k=j$. Nonnegativity handles $k\notin T$. Hence
\begin{equation}
 p^j\ge\bar a(\Dscr),\qquad \|p^j-\bar a(\Dscr)\|_1=s(\Dscr)\ge0.
 \label{eq:cover-distance}
\end{equation}
We use the following consequence of the property in Theorem~\ref{thm:family}(c). Whenever $p\in\Delta_T$
and $\Dscr\le_jC_{q^j}$ for $j\in T$,
\begin{equation}
 \mu(\Dscr\sd C_{p})\le\sum_{j\in T}\mu(C_{q^j}\sd C_{p}).
 \label{eq:face-comparison}
\end{equation}
Indeed, if $g\in \Dscr\setminus C_{p}$, some $g_j\le0$ for $j\in T$:
strictly positive payoffs throughout $T$ would force acceptance by $C_{p}$.
The corresponding family member must accept $g$, so it differs from $C_{p}$.
If $g\in C_{p}\setminus \Dscr$, some $g_j>0$ for $j\in T$, apart from the null
coordinate hyperplanes; otherwise the property in Theorem~\ref{thm:family}(c) forces rejection by $C_{p}$.
The corresponding family member must reject $g$, so again it differs from $C_{p}$.
This proves the bound, after discarding the finitely many null exceptions
to the background comparisons.

If $s(\Dscr)=0$, all the indices attaining the maxima equal $p=\bar a(\Dscr)\in\Delta_T$;
\eqref{eq:face-comparison} gives $\Dscr\approx C_{p}$. Conversely, (b) gives
$s(\Dscr)=0$ for such a model. If $T(\Dscr)$ is empty, $s(\Dscr)=1$.
\end{proof}

The adjustment must remain positive even when we compare a model with all
its possible background improvements. The next lemma controls the limiting
case: a sequence with vanishing allowance must approach a family member.

\begin{lemma}[Approaching a family member]\label{lem:stability}
If $s(\Dscr_k)\to0$, some subsequence and some $p\in\Delta$ satisfy
$\bar a(\Dscr_k)\to p$ and $\mu(\Dscr_k\sd C_{p})\to0$.
\end{lemma}
\begin{proof}
Pass to a subsequence with the nonempty set $T(\Dscr_k)=T$ fixed and
$\bar a(\Dscr_k)\to p$. The coordinate sums tend to one, so $p\in\Delta_T$.
Choose $p^{j,k}$ attaining the maximum in $a_j(\Dscr_k)$ for $j\in T$. By
\eqref{eq:cover-distance}, each $p^{j,k}\to p$.
Continuity in measure and \eqref{eq:face-comparison} give
\[
 \mu(\Dscr_k\sd C_{p})\le\sum_{j\in T}\mu(C_{p^{j,k}}\sd C_{p})\longrightarrow0.
\]
\end{proof}

To adjust the levels while preserving background monotonicity, put
\begin{equation}
 b_i(\Dscr)\coloneqq\inf_{\Escr\le_i\Dscr}\left(a_i(\Escr)+\frac{s(\Escr)}{2n}\right)
 \quad(i\in T(\Dscr)).
 \label{eq:envelope}
\end{equation}
Every background improvement $\Escr\le_i\Dscr$ also has an available index at $i$. Taking the
infimum over all such improvements makes the adjustment compatible with all
of them at once. The next lemma shows that the adjustment is positive
whenever it is needed.

\begin{lemma}[A positive adjustment that respects accuracy]\label{lem:reserve}
At a world with an available index,
\begin{equation}
 a_i(\Dscr)\le b_i(\Dscr)\le a_i(\Dscr)+\frac{s(\Dscr)}{2n},\qquad
 \sum_{i\in T(\Dscr)}b_i(\Dscr)\le1-\frac{s(\Dscr)}2.
 \label{eq:budget}
\end{equation}
If $\Dscr\le_i\Escr$ and $\Escr$ has an available index at $i$, then $b_i(\Dscr)\ge b_i(\Escr)$.
Moreover, $b_i(\Dscr)=a_i(\Dscr)$ exactly when $\Dscr\approx C_{p}$ for some $p$.
In particular, $a_i(C_{p})=b_i(C_{p})=p_i$.
\end{lemma}
\begin{proof}
Every term in the infimum is at least $a_i(\Dscr)$ by Lemma~\ref{lem:levels}(a),
and the term $\Escr=\Dscr$ gives the upper bound. Summing gives
$\sum_{i\in T(\Dscr)}b_i(\Dscr)\le1-s(\Dscr)+|T(\Dscr)|s(\Dscr)/(2n)\le1-s(\Dscr)/2$.
The set of background improvements of $\Dscr$ is contained in that of $\Escr$,
so its infimum is no smaller. This proves monotonicity.

If $\Dscr\approx C_{p}$, then $s(\Dscr)=0$ and the bounds give equality.
Conversely, suppose $b_i(\Dscr)=a_i(\Dscr)$. Choose $\Escr_k\le_i\Dscr$ with
\[
 a_i(\Escr_k)+s(\Escr_k)/(2n)\longrightarrow a_i(\Dscr).
\]
Since $a_i(\Escr_k)\ge a_i(\Dscr)$ and $s(\Escr_k)\ge0$, we have
$a_i(\Escr_k)\to a_i(\Dscr)$ and $s(\Escr_k)\to0$.
Lemma~\ref{lem:stability} gives a subsequence converging in measure to some
$C_{p}$, with $p_i=a_i(\Dscr)$. Passing to the limit in $\Escr_k\le_i\Dscr$ gives
$C_{p}\le_i\Dscr$: the measure of $E_i(C_{p})\setminus E_i(\Dscr)$ is at most
$\mu(C_{p}\sd \Escr_k)$. Choose $q$ attaining the maximum in $a_i(\Dscr)$. Then
\[
 C_{p}\le_i\Dscr\le_iC_{q},\qquad p_i=q_i=a_i(\Dscr).
\]
Lemma~\ref{lem:levels}(b) gives $p=q$, and the sandwich gives $\Dscr\approx C_{p}$.
\end{proof}

We can now define the ranking. Compare the pairs
$(a_i(\Dscr),b_i(\Dscr))$ coordinate by coordinate, with larger values better;
assign $(-1,-1)$ when no index is available. If two models have the same pair,
use their background comparison, except that all models with pair $(t,t)$,
$t\ge0$, tie. Such pairs belong only to family members, by
Lemma~\ref{lem:reserve}.

Denote the resulting relations by $\preceq_i$, with strict parts $\prec_i$.
Changes on null sets do not affect the ranking. The next formulas show exactly
how each selected family member compares with an arbitrary acceptance set;
they will give the closedness and competitiveness required by (S).
Since a family member has pair $(p_i,p_i)$, its comparisons are exactly
\begin{equation}
 C_{p}\preceq_i\Dscr\quad\Longleftrightarrow\quad
 \begin{cases}p_i\ge b_i(\Dscr),&i\in T(\Dscr),\\
 \text{always},&i\notin T(\Dscr),\end{cases}\label{eq:outgoing}
\end{equation}
\begin{equation}
 \Dscr\preceq_iC_{p}\quad\Longleftrightarrow\quad
 i\in T(\Dscr)\ \text{and}\ a_i(\Dscr)\ge p_i.\label{eq:incoming}
\end{equation}
In particular, the sets of indices giving a comparison are closed:
\begin{equation}
 \{p:C_{p}\preceq_i\Dscr\}=\{p:p_i\ge b_i(\Dscr)\}\quad(i\in T(\Dscr)),\label{eq:sections}
\end{equation}
and that set is all of $\Delta$ at a world with no available index.

We have built a family through every credal set $M$. The next lemma ensures
that this is enough: every coherent acceptance set agrees, up to a null set,
with $\Dscr_M$ for some $M$.
\begin{lemma}\label{lem:coherent}
For any coherent set of desirable gambles $\Dscr$, the set
\[
 M=\{m\in\Delta:m\cdot g\ge0\text{ for all }g\in \Dscr\}
\]
is a credal set, and
$\{\lpr{M}>0\}\subseteq \Dscr\subseteq\{\lpr{M}\ge0\}$.
Consequently $\Dscr$ is Lebesgue measurable and $\Dscr\approx \Dscr_{M}$.
\end{lemma}

\begin{proof}
The set $M$ is closed and convex as an intersection of closed half-spaces with
$\Delta$. Let $C=\overline \Dscr$. It is a closed convex cone and contains zero,
since $te_1\in \Dscr$ for $t>0$. No strictly negative vector lies in $C$: $\Dscr$
contains no nonzero $g\le0$, since otherwise $-g\in \Dscr$ and addition would
put zero in $\Dscr$; a strictly negative vector has a neighbourhood of such
excluded points.

We first show that if $h\notin C$, some $m\in M$ has $m\cdot h<0$.
Choose $c\in C$ nearest to $h$ in Euclidean distance, which exists because
$C$ is nonempty and closed, and put $w=c-h\ne0$. For $x\in C$, minimality
along the segment from $c$ to $x$ gives
\[
 \|w+t(x-c)\|^2\ge\|w\|^2\quad(0<t\le1).
\]
Dividing the expanded inequality by $t$ and letting $t\to0$ yields
$w\cdot(x-c)\ge0$. Taking $x=0$ and $x=2c$ shows $w\cdot c=0$.
Hence $w\cdot x\ge0$ on $C$, while $w\cdot h=-\|w\|^2<0$.
Since every $e_i\in \Dscr$, we have $w_i\ge0$, and normalization gives
$m=w/\sum_iw_i\in M$ with $m\cdot h<0$.

Applying this to $h=-\one\notin C$ proves $M\ne\varnothing$.
The inclusion $\Dscr\subseteq\{\lpr{M}\ge0\}$ follows from the definition of
$M$. For the other inclusion, suppose $\lpr{M}(g)=2\varepsilon>0$ and put
$h=g-\varepsilon\one$. Then $\lpr{M}(h)=\varepsilon>0$, so the separation
just proved implies $h\in C$. Choose $h'\in \Dscr$ with
$\|h'-h\|_\infty<\varepsilon$. Now $g-h'$ has every coordinate positive,
so $g-h'\in \Dscr$ and $g=h'+(g-h')\in \Dscr$.

Both $\Dscr$ and $\Dscr_{M}$ thus lie between $\{\lpr{M}>0\}$ and
$\{\lpr{M}\ge0\}$. Their difference from the first, open set is contained
in the null zero set of $\lpr{M}$. This proves measurability and
$\Dscr\approx \Dscr_{M}$.
\end{proof}

We now put the construction together. The theorem verifies the background
accuracy requirements, keeps the chosen family undominated, and shows that
allowing the target to vary covers every coherent acceptance set.

\begin{theorem}[Full coverage by accuracy rankings]\label{thm:coverage}
For each target $M$, the relations above are preorders on all measurable
models. They respect both weak and strict background comparisons:
\begin{equation}
 \boxed{\quad \Dscr\le_i\Escr\ \Longrightarrow\ \Dscr\preceq_i\Escr,
 \qquad \Dscr<_i\Escr\ \Longrightarrow\ \Dscr\prec_i\Escr.\quad}\label{eq:accuracy}
\end{equation}
Their undominated models are exactly the selected family modulo null sets.
Each $C_{p}$ is competitive on $\supp p$, the sets of comparison indices are closed,
and every model outside the family modulo null sets is strictly dominated at
every world by a family member.

Letting $M$ range over all credal sets gives a class of rankings satisfying
\textbf{(E)}, \textbf{(LA)}, \textbf{(A)}, and \textbf{(S)}. With coherence
understood modulo null sets, they also satisfy \textbf{(C)}.
\end{theorem}
\begin{proof}
The relations are reflexive. Along a comparison chain the comparison pairs
can only decrease coordinatewise; if the first and last pairs agree, the
whole chain lies in one group, where comparison is either background
inclusion or a tie between every pair of models. Thus the relations are transitive.

For accuracy, $\Dscr\le_i\Escr$ gives $\mathcal P_i(\Escr)\subseteq\mathcal P_i(\Dscr)$,
so $a_i(\Dscr)\ge a_i(\Escr)$ when $\Escr$ has an available index. The set of background improvements
of $\Dscr$ is contained in that of $\Escr$, so $b_i(\Dscr)\ge b_i(\Escr)$ as well.
At a fixed world, a model with an available index also ranks above one without. Unequal comparison pairs
therefore make a strict background improvement strict; equal pairs retain
background inclusion. The only possible exception is a pair $(t,t)$ with $t\ge0$, where $\Dscr\approx C_{p}$, $\Escr\approx C_{q}$ and $p_i=q_i=t$.
But $C_{p}\le_iC_{q}$ then forces $p=q$ by Lemma~\ref{lem:levels}(b): (CC) gives
$p_j\le q_j$ for every $j\ne i$, and the coordinate sums are equal.
Thus $\Dscr\approx \Escr$, ruling out a strict background comparison in this
exceptional case. This proves \eqref{eq:accuracy}.

For competitiveness, any $i\in\supp p$ with no available index supplies the comparison
by \eqref{eq:outgoing}. Otherwise, failure at every world in $\supp p$ would
require $p_i<b_i(\Dscr)$ there, contradicting
$1=\sum_{i\in\supp p}p_i<\sum_{i\in\supp p}b_i(\Dscr)\le1$.
Closedness follows from \eqref{eq:sections}.
If $\Dscr\preceq_iC_{p}$ at every world, \eqref{eq:incoming} and the sum bound force
$a_i(\Dscr)=p_i$ for all $i$ and $s(\Dscr)=0$, hence $\Dscr\approx C_{p}$.
So family members are undominated.

For a model outside the selected family up to null sets, set $\beta_i=b_i(\Dscr)$ at worlds with an available index and zero
elsewhere. By \eqref{eq:budget}, $\sum_i\beta_i<1$. The index
$q=\beta+(1-\sum_i\beta_i)\one/n$ has $q_i>b_i(\Dscr)\ge a_i(\Dscr)$ at worlds with an available index, and its pair exceeds $(-1,-1)$ at worlds with no available index. Thus $C_{q}\prec_i\Dscr$
everywhere.

All models are evaluated, giving (E); the selected coherent undominated models
give (LA); and $p\mapsto C_{p}$ gives (S). The characterization of undominated
models and the fact that changes on null sets do not affect the ranking give (C).
Finally, any coherent $\Dscr$ agrees a.e. with $\Dscr_{M}$ by Lemma~\ref{lem:coherent}.
The family for that target has $K_r=M$, so $\Dscr$ ties everywhere with its
undominated member $C_{r}$. Thus $\Dscr$ is undominated under this ranking,
establishing (A), also for all models coherent modulo null sets.
\end{proof}

\begin{remark}[Uniform target index]
The target can be selected at the uniform index without changing the ranking.
Put $\gamma(p)=n\min_i p_i$ and $f(p)=(1-\gamma(p))p+\gamma(p)r$.
Then $f$ is continuous, $f(\one/n)=r$, and
$\supp f(p)\subseteq\supp p$: on the boundary $f(p)=p$, while interior
indices have full support. Hence $\sigma_p=\Dscr_{K_{f(p)}}$ is competitive on
$\supp p$, and its sets of comparison indices are closed as continuous inverse
images. Thus $\sigma_{\one/n}=\Dscr_{M}$. This reparametrization is applied after
constructing the ranking; it need not preserve (CC) in the new indices.
\end{remark}

\section{Constructing a bounded real-valued loss}

The same family also gives a numerical loss. Its dependence on the target is
through the family $K_p$ and the available indices $\mathcal P_i(\Dscr)$.
We will make every $C_p=\Dscr_{K_p}$ have loss $1-p_i$ at $\omega_i$;
since $C_r=\Dscr_M$, this includes the target.

This requires quantitative bounds on how the family changes with its index.
Theorem~\ref{thm:family} supplies
\[
 \dH(K_p,K_q)\le A\|p-q\|_1,\qquad \upprob{K_p}(\omega_i)\le Bp_i,
 \qquad B=1/\lambda.
\]
An explicit choice is $A=1+2/\delta$, where $\delta=\min_{i:d_i>0}d_i$:
the minimum defining $\theta$ is $1/\delta$-Lipschitz, so
\eqref{eq:hausdorff} applies. For singleton targets take $A=1$.

Fix $\mu$ to be standard Gaussian measure, put $\kappa=8n^2(A+2B)$, and choose
$0<\eta<1/(n\kappa)$. Using the same sets of available indices, define
\begin{equation}
 L_\eta(\Dscr,\omega_i)\coloneqq
 \begin{cases}
 \displaystyle\min_{p\in\mathcal P_i(\Dscr)}
 \left[1-p_i-\eta\mu\big(E_i(\Dscr_{K_p})\setminus E_i(\Dscr)\big)\right],
       &\mathcal P_i(\Dscr)\ne\varnothing,\\[3pt]
 2+\mu(E_i(\Dscr)),&\mathcal P_i(\Dscr)=\varnothing.
 \end{cases}\label{eq:finite-loss}
\end{equation}
The minimum exists: the set of available indices is compact, and continuity in measure of
$\Dscr_{K_p}$ makes the objective continuous in $p$. The correction rewards errors removed from a family member used for comparison; the bound on $\eta$ keeps these rewards within the sum bound needed for competitiveness.

The next two lemmas bound the size of the error reductions. They let us
choose the reward $\eta$ small enough that no model can beat a family member
at every world. Both use elementary estimates for Gaussian measure. Here a random vector with
law $\mu$ has independent standard normal coordinates; $\|\cdot\|_2$ denotes
the Euclidean norm.

\begin{lemma}[Gambles with lower prevision near zero]\label{lem:thin}
For a credal set $K$ and $\delta\ge0$,
\[
 \mu\{g:|\lpr{K}(g)|\le\delta\|g\|_\infty\}\le4n\delta.
\]
Consequently $\mu(\Dscr_{K}\sd \Dscr_{H})\le4n\,\dH(K,H)$ for any credal sets $K,H$.
\end{lemma}
\begin{proof}
For $\delta>1/2$ the bound is automatic. Otherwise write
$g=h+t\one$, where $t=n^{-1}\sum_i g_i$ and $\sum_i h_i=0$.
Equation~\eqref{eq:constant} and $|\lpr{K}(h)|\le\|h\|_\infty$ give
\[
 \|g\|_\infty\le2\|h\|_\infty+|t+\lpr{K}(h)|.
\]
Thus $|\lpr{K}(g)|\le\delta\|g\|_\infty$ implies
$|t+\lpr{K}(h)|\le4\delta\|h\|_\infty$.
For standard Gaussian $g$, the density factorizes in these coordinates
because $\|g\|_2^2=\|h\|_2^2+nt^2$. Thus $t$ and $h$ are independent,
$t$ has variance $1/n$, and its density is bounded by
$\sqrt{n/(2\pi)}$. Conditional on $h$, the displayed condition restricts $t$
to an interval of length $8\delta\|h\|_\infty$. Since
\[
 \mathbb E\|h\|_\infty\le\mathbb E\|h\|_2
 \le\sqrt{\mathbb E\|h\|_2^2}=\sqrt{n-1},
\]
its probability is at most
$8\delta\sqrt{n(n-1)/(2\pi)}\le4n\delta$.
Finally, if $\Dscr_{K}$ and $\Dscr_{H}$ disagree, their lower previsions have opposite
signs, up to null sets. By \eqref{eq:phi-estimates}, this implies
$|\lpr{K}(g)|\le\dH(K,H)\|g\|_\infty$, giving the consequence.
\end{proof}

\begin{lemma}[Errors when outside probability is small]\label{lem:face-errors}
For nonempty $T\subseteq[n]$, let
$A_T^+=\{g:g_i>0\ (i\in T)\}$ and $A_T^-=\{g:g_i\le0\ (i\in T)\}$.
If $\upprob{K_p}(\omega_i)\le Bp_i$ for all $i$, then
\[
 \mu(A_T^+\setminus \Dscr_{K_p}),\quad\mu(\Dscr_{K_p}\cap A_T^-)
 \ \le\ 8n^2B\,p(T^c).
\]
\end{lemma}
\begin{proof}
Put $\varepsilon=Bp(T^c)$. If $\varepsilon>1/2$ the bounds are automatic.
Otherwise every $m\in K_p$ has $m(T^c)\le\varepsilon$ and
$m(T)\ge1-\varepsilon$.
For $g\in A_T^+\setminus \Dscr_{K_p}$, choose $m\in K_p$ with $m\cdot g\le0$
and put $\gamma=\min_{i\in T}g_i$. Then
\[
 0\ge m\cdot g\ge(1-\varepsilon)\gamma-\varepsilon\|g\|_\infty,
\]
so $\gamma\le2\varepsilon\|g\|_\infty$.
For $g\in \Dscr_{K_p}\cap A_T^-$, discard the null coordinate hyperplanes and
put $\gamma=\min_{i\in T}|g_i|$. All coordinates on $T$ are negative, while
$\lpr{K_p}(g)>0$, so any $m\in K_p$ gives
\[
 0<m\cdot g\le-(1-\varepsilon)\gamma+\varepsilon\|g\|_\infty.
\]
Again $\gamma\le2\varepsilon\|g\|_\infty$. Both error regions therefore lie,
up to null sets, in
$\bigcup_{i\in T}\{g:|g_i|\le2\varepsilon\|g\|_\infty\}$.
Apply Lemma~\ref{lem:thin} with $K=\{e_i\}$ to bound each term by
$8n\varepsilon$, and sum over at most $n$ coordinates.
\end{proof}

With these estimates, the numerical proof follows the same pattern as the
ordinal one: establish a sum bound, show that family members attain it,
and use it to rule out domination of the target.

\begin{theorem}[Full coverage by bounded losses]\label{thm:numerical}
For each target $M$, the loss \eqref{eq:finite-loss} takes values in
$[0,1]\cup[2,3]$ and evaluates every measurable model. It satisfies
\[
 \Dscr\le_i\Escr\ \Longrightarrow\ L_\eta(\Dscr,\omega_i)\le L_\eta(\Escr,\omega_i),\qquad
 \Dscr<_i\Escr\ \Longrightarrow\ L_\eta(\Dscr,\omega_i)<L_\eta(\Escr,\omega_i).
\]
Moreover:
\begin{enumerate}[label=(\alph*)]
\item $L_\eta(\Dscr_{K_p},\omega_i)=1-p_i$;
\item the selected models are undominated and competitive on their index
supports, and all sets of indices giving a comparison are closed;
\item every model outside the selected family modulo null sets is strictly
dominated at every world by a selected model.
\end{enumerate}
Letting $M$ range over all credal sets gives \textbf{(E)}, \textbf{(LA)},
\textbf{(A)}, and \textbf{(S)}, as well as \textbf{(C)} with coherence
understood modulo null sets.
\end{theorem}
\begin{proof}
We give the sum bound argument, writing $C_{p}=\Dscr_{K_p}$. For
$T=T(\Dscr)\ne\varnothing$, choose available indices $p^i\in\mathcal P_i(\Dscr)$ and put
\[
 \rho=1-\sum_{i\in T}p^i_i,\qquad
 h_i=\mu\big(E_i(C_{p^i})\setminus E_i(\Dscr)\big).
\]
Let $\alpha_k=p^k_k$ for $k\in T$ and $\alpha_k=0$ otherwise. By (CC),
$p^j\ge\alpha$ for every $j\in T$, so $\|p^j-\alpha\|_1=\rho$. Thus
\[
 \rho\ge0,\qquad p^j(T^c)\le\rho,\qquad
 \|p^i-p^j\|_1\le2\rho.
\]
Lemmas~\ref{lem:thin} and~\ref{lem:face-errors} give
$\mu(C_{p}\sd C_{q})\le4nA\|p-q\|_1$ and bound each of the two error regions by
$8n^2B\,p(T^c)$.
An error removed from $C_{p^j}$ by $\Dscr$ either witnesses disagreement with
another family member used for comparison or lies in one of these two error regions. Indeed,
if $g\in \Dscr\setminus C_{p^j}$ and some $g_i\le0$ for $i\in T$, the family member indexed by $p^i$ must accept $g$; otherwise $g\in A_T^+\setminus C_{p^j}$. If
$g\in C_{p^j}\setminus \Dscr$ and some $g_i>0$ for $i\in T$, that family member must reject $g$; otherwise $g\in C_{p^j}\cap A_T^-$. These statements hold outside
the null exceptional sets of the background comparisons. Hence
\[
 h_j\le\sum_{i\in T\setminus\{j\}}\mu(C_{p^i}\sd C_{p^j})
       +16n^2B\,p^j(T^c)\le\kappa\rho,
\]
and therefore
\begin{equation}
 \sum_{i\in T}(p^i_i+\eta h_i)
 \le1-(1-n\kappa\eta)\rho\le1.\label{eq:numerical-budget}
\end{equation}
Choose indices attaining the loss minimum and write $c_i(\Dscr)=1-L_\eta(\Dscr,\omega_i)$
for $i\in T$. These numbers are nonnegative and sum to at most one, proving
the loss range. Equality forces $\rho=0$, all chosen indices equal to some
$p$, and every $h_i=0$. Thus $E_i(\Dscr)=E_i(C_{p})$ a.e. at any world with an available index,
and $\Dscr\approx C_{p}$.

For a family member, its own index is available at every world, so
$c_i(C_{p})\ge p_i$; the sum bound forces equality, proving (a). If $\Dscr$ weakly
dominates $C_{p}$, it has an available index at every world and $c_i(\Dscr)\ge p_i$ for all $i$.
Equality in the sum bound gives $\Dscr\approx C_{q}$ for some $q$, and then
$q_i=c_i(\Dscr)\ge p_i$ for all $i$, forcing $q=p$. Hence $C_{p}$ is undominated.

For accuracy, if $\Dscr\le_i\Escr$ and $\Escr$ has an available index, every index available for $\Escr$ is available for $\Dscr$
and the correction increases by $\eta\mu(E_i(\Escr)\setminus E_i(\Dscr))$. Thus
\[
 L_\eta(\Dscr,\omega_i)\le L_\eta(\Escr,\omega_i)-\eta\mu(E_i(\Escr)\setminus E_i(\Dscr)).
\]
If both models have no available index, monotonicity and strictness follow from the
error measure; if only $\Dscr$ has an available index, they follow from the disjoint loss ranges.

Competitiveness follows from the sum bound: strict improvement over $C_{p}$ at
every world in $\supp p$ would require $\sum_{i\in\supp p}c_i(\Dscr)>1$.
The set of comparison indices is $\{p:p_i\ge1-L_\eta(\Dscr,\omega_i)\}$, which is closed.
Finally, set $\beta_i=c_i(\Dscr)$ at worlds with an available index and zero elsewhere.
For a model outside the selected family up to null sets $\sum_i\beta_i<1$, so
$q=\beta+(1-\sum_i\beta_i)\one/n$ supplies strict dominance everywhere.
The axioms follow as in Theorem~\ref{thm:coverage}, with full coverage again
supplied by $K_r=M$.
\end{proof}

\paragraph{Why the target is undominated.}
The proof gives $\sum_i L_\eta(\Dscr,\omega_i)\ge n-1$ for every acceptance
set, while $L_\eta(C_p,\omega_i)=1-p_i$ for every family member.
In particular $C_r=\Dscr_M$ attains total loss $n-1$.
A model dominating it would have a smaller total, contradicting the bound.

This loss gives another ranking on the same family; it does not represent the
preceding ranking, which has incomparabilities, and need not preserve all its
strict comparisons. The construction imposes no additivity requirement on the loss.

\appendix
\section{Why upper probabilities alone are too restrictive}\label{app:ustar}

This appendix is not needed for either construction. It explains why an earlier
condition, which used only upper probabilities, excluded targets that the
present proof can include.

An earlier sufficient condition for (CC) was
\begin{equation*}
 \upprob{K_p}(\omega_i)\le\upprob{K_q}(\omega_i),\quad
 \upprob{K_q}(\omega_j)\le\upprob{K_p}(\omega_j),\quad i\ne j
 \quad\Longrightarrow\quad p_i\le q_i.
 \tag{U$^*$}\label{eq:ustar}
\end{equation*}
Unlike (CC), it requires an index comparison even when the upper-probability
inequalities do not arise from a single acceptance set compared at two worlds.

For example, the vacuous family \eqref{eq:vacuous} in three worlds gives,
at $p=(.70,.20,.10)$ and $q=(.65,.13,.22)$, upper probability vectors
$(.90,.40,.30)$ and $(.91,.39,.48)$. The first coordinate increases and the
second decreases, but $p_1>q_1$, violating (U$^*$). This causes no failure
of (CC): an acceptance set $\Dscr$ with $\Dscr\le_1\Dscr_{K_p}$ and
$\Dscr\le_2\Dscr_{K_q}$ would also require
$\lowprob{K_p}(\omega_1)\le\lowprob{K_q}(\omega_1)$, whereas these values are $.60$ and $.52$.

\begin{proposition}\label{prop:ustar}
Suppose a family is continuous in Hausdorff distance and satisfies both
$K_p\subseteq\Delta_{\supp p}$ and (U$^*$).
If $K_p$ contains a simplex vertex, then $|\supp p|\le2$. Consequently such
families cannot contain a credal set containing a vertex whose distributions collectively assign positive probability
to at least three worlds; in particular, they cannot contain $\Delta$ for
$n\ge3$.
\end{proposition}
\begin{proof}
Suppose $e_i\in K_p$ and choose distinct $j,k\in\supp p\setminus\{i\}$.
Here $i\in\supp p$ because $K_p\subseteq\Delta_{\supp p}$. Hold all other coordinates fixed and
vary $i,j,k$. With $c=p_i+p_j+p_k$, the interior of this triangular slice is
parametrized by
\[
 V=\{(s,t):s>0,\ t>0,\ s+t<c\},\qquad
 x_i=s,\quad x_j=t,\quad x_k=c-s-t.
\]
The map $F(s,t)=(\upprob{K_x}(\omega_i),\upprob{K_x}(\omega_j))$ is continuous. It is injective:
equality of both outputs, applying (U$^*$) in both directions and exchanging
$i,j$, gives equality of both $x_i,x_j$.

Brouwer's invariance of domain theorem says that a continuous injective map
from an open subset of $\R^2$ into $\R^2$ has open image. Thus $F(V)$ is open.
It contains $F(p_i,p_j)$, whose first coordinate is one because $e_i\in K_p$.
An open neighbourhood of this point includes first coordinates greater than
one, impossible for upper probabilities. Hence $|\supp p|\le2$.
\end{proof}

This restriction concerns the stronger condition (U$^*$).
The coverage construction needs only (CC) and does not use invariance of domain.

\end{document}
